\begin{afterword}
\summaryhead

The author describes growing up a technophile. At nine the author read Jerry Pournelle's \textsc{A Step Farther Out}, a reply to Ehrlich, the Club of Rome and the opponents of nuclear power, and took from it a picture of history: scientists and engineers on one side, the opponents of vaccination and steam engines on the other, and in the middle the ``pretenders to Deep Wisdom'' who call for regulation and committees.

The adult author still thinks this picture was not ``too wrong'' about who the Bad Guys were, and cites an estimate that slow drug approval in the United States causes far more deaths than it prevents. The mistake was to think that avoiding everything the Bad Guys did makes one a Good Guy, so that any caution felt like joining the middle. In a debate about whether nanotechnology would favour attack over defense, which the author dates, with ``I think,'' to 1997 or 1998, the young author saw for the first time that extinction was possible. Yet little changed: the plan to race to AI stayed, now justified as a way to get there before nanotechnology. Even a sharp Traditional Rationalist, the author concludes, falls short of what is needed.

\responsehead

As memoir the post is accurate, and the archive confirms it. The Extropians mailing list contains the 1997 debate on nanotechnology, with the young author's argument from nuclear weapons; the 1996 slogan ``Get to the Singularity as fast as possible!''; and a 1999 statement that the hurry was about humanity wiping itself out. So the post's central confession is on record: the author kept the old plan of racing to AI and gave it a new reason.

The post also defends, in the present tense, the childhood habit of being ``allergic'' to anyone who says that technology has risks as well as benefits. The author says there is ``a historical record showing over-conservativeness.'' The evidence given is one estimate about one regulator, the US Food and Drug Administration, and a list of old disputes over vaccination, anaesthesia, steam engines and heliocentrism. The post does not argue that what holds for drug approval holds for technology in general. It also appears to misreport the estimate: the figures seem to be Dale Gieringer's, which count deaths per decade, not per year. Both figures are off by the same factor, so the comparison between them still holds.

The drug estimate cannot address the risk that the rest of the post is about. The estimate compares lives lost to slow approval with lives lost to harmful drugs. The young author's mistake, in the post's own words, was to act as if the Future could not be lost: ``Lives could be lost, but not the Future.'' A record of drug approvals says nothing about risks to the whole future. The post keeps both the old allergy and the new lesson, and does not say how to tell which kind of risk one is facing.

The post dismisses the middle position with a label instead of an argument. It calls those who hold it ``pretenders to Deep Wisdom,'' and says they claim ``in defiance of brute historical fact'' that science of itself is ``neither good nor evil.'' But a record of good uses of science does not contradict a claim about what science is in itself. From this one set of opponents the post draws a general rule: seeking a middle way ``is usually wrong.''

In short: the memoir is accurate and the archive confirms it, but the post's defence of the author's allergy to talk of risk rests on a single estimate about drug approval, apparently misreported by a factor of ten, which says nothing about the risk to the whole future that the post is about.
\end{afterword}
